%% kapa: exam=1 round=19 year=2008 period=1 form=A subject=civil_law subject_ko=민법 seq=9 number=9
%% source: Q-Net 감정평가사 기출문제 (cst00309) artlSeq=1001962 file=extracted/1008/1001962/2008년 제19회 감정평가사 1차시험 기출문제/1교시.hwp
%% conversion: kordoc 4.24.51 -> tools/gen1.py (auto); review=proofread-sonnet; flags=none
\begin{question}{9}
착오에 의한 의사표시에 관한 설명으로 틀린 것은?(다툼이 있는 경우에는 판례에 의함)
\begin{choices}
  \item 화해의 목적인 분쟁 이외의 사항에 착오가 있는 때에는 착오를 이유로 화해계약을 취소할 수 있다.
  \item 법률에 관한 착오라고 하더라도 그것이 법률행위 내용의 중요부분에 관한 것인 때에는 표의자는 그 의사표시를 취소할 수 있다.
  \item 법률행위 내용의 중요부분에 착오가 있더라도 표의자에게 중대한 과실이 있으면 착오를 이유로 의사표시를 취소할 수 없다.
  \item 부동산 매매에서 시가에 관한 착오는 원칙적으로 법률행위 내용의 중요부분에 관한 착오에 해당한다.
  \item 법률행위 내용의 중요부분에 관한 착오란 표의자에게 그러한 착오가 없었더라면 그 의사표시를 하지 않으리라고 생각될 정도로 중요한 것이어야 하고, 보통 일반인도 표의자의 처지에 섰더라면 그러한 의사표시를 하지 않았으리라고 생각될 정도로 중요한 것이어야 한다.
\end{choices}
\end{question}
\answer{4}
